% =====================================================================
%  ATTPCROOTv2-OpenKF
%  Track Reconstruction with the UKF and genfit Fitters
%  A user manual for the full simulation and experimental pipelines
% =====================================================================
\documentclass[11pt,a4paper]{article}

\usepackage[utf8]{inputenc}
\usepackage[T1]{fontenc}
\usepackage{lmodern}
\usepackage[margin=2.4cm]{geometry}
\usepackage{microtype}
\usepackage{graphicx}
\usepackage{xcolor}
\usepackage{booktabs}
\usepackage{longtable}
\usepackage{array}
\usepackage{enumitem}
\usepackage{amsmath}
\usepackage{listings}
\usepackage{tcolorbox}
\usepackage{fancyhdr}
\usepackage[hidelinks,colorlinks=true,linkcolor=teal!70!black,urlcolor=teal!70!black]{hyperref}

% ---------- colors ----------
\definecolor{codebg}{HTML}{F6F8FA}
\definecolor{codeframe}{HTML}{D0D7DE}
\definecolor{kw}{HTML}{0550AE}
\definecolor{cmt}{HTML}{6E7781}
\definecolor{str}{HTML}{0A7D33}
\definecolor{accent}{HTML}{1F6F6F}

% ---------- listings ----------
\lstdefinestyle{cpp}{
  language=C++,
  backgroundcolor=\color{codebg},
  basicstyle=\ttfamily\footnotesize,
  keywordstyle=\color{kw}\bfseries,
  commentstyle=\color{cmt}\itshape,
  stringstyle=\color{str},
  frame=single,
  rulecolor=\color{codeframe},
  breaklines=true,
  showstringspaces=false,
  columns=fullflexible,
  numbers=none,
  morekeywords={Float_t,Int_t,Double_t,Bool_t,TString,kTRUE,kFALSE,auto,nullptr,override,unique_ptr,make_unique,shared_ptr}
}
\lstdefinestyle{bash}{
  language=bash,
  backgroundcolor=\color{codebg},
  basicstyle=\ttfamily\footnotesize,
  keywordstyle=\color{kw},
  commentstyle=\color{cmt}\itshape,
  stringstyle=\color{str},
  frame=single,
  rulecolor=\color{codeframe},
  breaklines=true,
  showstringspaces=false,
  columns=fullflexible,
  numbers=none
}
\lstset{style=cpp}

% ---------- callout boxes ----------
\newtcolorbox{notebox}{colback=teal!4,colframe=accent,boxrule=0.6pt,arc=2pt,
  left=6pt,right=6pt,top=4pt,bottom=4pt,fonttitle=\bfseries,title=Note}
\newtcolorbox{warnbox}{colback=red!4,colframe=red!55!black,boxrule=0.6pt,arc=2pt,
  left=6pt,right=6pt,top=4pt,bottom=4pt,fonttitle=\bfseries,title=Important}

% ---------- headers ----------
\pagestyle{fancy}
\fancyhf{}
\renewcommand{\headrulewidth}{0.3pt}
\lhead{\small ATTPCROOTv2-OpenKF}
\rhead{\small UKF / genfit Pipeline Manual}
\cfoot{\thepage}

\newcommand{\code}[1]{\texttt{\small #1}}
\newcommand{\file}[1]{\texttt{\small #1}}

% =====================================================================
\begin{document}

% ---------------- title ----------------
\begin{titlepage}
\centering
\vspace*{2.5cm}
{\Huge\bfseries Track Reconstruction in ATTPCROOTv2-OpenKF\par}
\vspace{0.6cm}
{\LARGE The UKF and genfit Fitting Pipelines\par}
\vspace{1.2cm}
{\large A User Manual\par}
\vspace{2.0cm}
{\large\itshape From simulation / raw HDF5 to fitted kinematics\par}
\vfill
{\large Active Target Time Projection Chamber (AT-TPC)\par}
\vspace{0.4cm}
{\normalsize Repository: \file{ATTPCROOTv2-OpenKF} \quad branch \file{OpenKF-Claude}\par}
\vspace{0.4cm}
{\normalsize \today\par}
\end{titlepage}

\tableofcontents
\newpage

% =====================================================================
\section{Introduction}
% =====================================================================

This manual describes how to reconstruct charged-particle tracks in the AT-TPC
using \textbf{ATTPCROOTv2-OpenKF}, covering the two track fitters available in the
framework:

\begin{itemize}[nosep]
  \item \textbf{UKF} --- the native \emph{Unscented Kalman Filter} built on the
  vendored OpenKF library (class \code{EventFit::AtFitterUKF}). It uses a CATIMA
  energy-loss model, sigma-point propagation through the magnetic field, and an
  RTS smoother with back-extrapolation to the reaction vertex.
  \item \textbf{genfit} --- the established GENFIT2 Kalman reference-track fitter
  (class \code{AtFITTER::AtGenfit}, \code{KalmanFitterRefTrack}), with SRIM-style
  energy-loss tables.
\end{itemize}

Both fitters consume the \emph{same} pattern-recognition output
(\code{AtPatternEvent}) and produce fitted tracks (\code{AtTrackingEvent} /
\code{AtFittedTrack}). They are interchangeable at the level of a single FairRoot
task, which makes it easy to fit identical track candidates with both methods and
compare.

\begin{notebox}
\textbf{When to use which.} The UKF is faster, native, easy to iterate, and uses
CATIMA energy loss --- prefer it for quick parameter scans and for setups where
its tuning has been validated. genfit uses analytic-Jacobian transport and tends
to reach better intrinsic resolution; prefer it when resolution is the priority.
The two are best used \emph{together}: reconstruct once, then fit with both.
\end{notebox}

The document is organized as a pipeline:

\begin{center}
\fbox{\parbox{0.95\textwidth}{\centering
\small\ttfamily
Simulation OR raw HDF5 \;$\rightarrow$\; Digitization / Unpacking \;$\rightarrow$\;
PSA (hits) \;$\rightarrow$\; [Space-charge corr.] \;$\rightarrow$\; PRA (track
candidates) \;$\rightarrow$\; \textbf{UKF / genfit fit} \;$\rightarrow$\;
Kinematics / Ex spectrum
}}
\end{center}

Section~\ref{sec:concepts} explains the stages.
Section~\ref{sec:env} covers the environment.
Section~\ref{sec:sim} walks the \emph{simulation} pipeline end-to-end
(\code{16C\_pp} example).
Section~\ref{sec:exp} walks the \emph{experimental} pipeline (\code{a1975}
example), including the two fitters side-by-side.
Section~\ref{sec:fitterapi} is the fitter API reference.
Section~\ref{sec:conventions} documents the crucial magnetic-field-sign /
handedness convention.
Sections~\ref{sec:datamodel}--\ref{sec:trouble} cover the output data model,
quick-reference cards, and troubleshooting.

% =====================================================================
\section{Pipeline concepts}\label{sec:concepts}
% =====================================================================

Every reconstruction is a chain of FairRoot tasks (\code{FairTask}) added to a
\code{FairRunAna}. Each task reads one or more ROOT branches and writes another.
The stages, in order:

\begin{description}[leftmargin=1.4cm,style=nextline]
  \item[Simulation / Unpacking] For simulation, a Geant4 (\code{TGeant4}) run
  transports the beam and reaction products and writes Monte-Carlo points. For
  experiment, \code{AtUnpackTask} (with \code{AtHDFUnpacker}) reads raw GET
  electronics traces from HDF5 into \code{AtRawEvent}.
  \item[Digitization] (Simulation only.) \code{AtClusterizeTask} converts MC
  ionization to electron clusters and \code{AtPulseTask} produces realistic pad
  signals (\code{AtRawEvent}), so simulation and experiment share the same
  downstream code.
  \item[PSA] \code{AtPSAtask} (pulse-shape analysis, typically \code{AtPSAMax})
  turns pad waveforms into 3D hits with charge: branch \code{AtEventH}.
  \item[Space-charge correction] (Experiment, optional.)
  \code{AtSpaceChargeCorrectionTask} applies distortion look-up tables and writes
  \code{AtEventCorrected}.
  \item[Pattern recognition (PRA)] \code{AtPRAtask} (Hough/triplet RANSAC) or
  \code{AtSampleConsensusTask} groups hits into clustered track candidates:
  branch \code{AtPatternEvent}.
  \item[Fitting] \code{AtFitterTask} (UKF or multi-hypothesis UKF) or
  \code{AtFitterTaskOld} (genfit) fits each candidate, producing
  \code{AtTrackingEvent} (a collection of \code{AtFittedTrack}) and optionally
  \code{AtFitMetadata}.
\end{description}

\begin{notebox}
The fitter only ever sees \code{AtPatternEvent}. This is why the fitting step can
be \emph{iterated independently}: reconstruct once up to PRA, then refit many
times with different fitter settings without re-running the slow stages.
\end{notebox}

% =====================================================================
\section{Environment and build}\label{sec:env}
% =====================================================================

\subsection{Sourcing the environment}

From the repository root, always source the environment before running any macro:

\begin{lstlisting}[style=bash]
cd /path/to/ATTPCROOTv2-OpenKF
source build/config.sh        # sets VMCWORKDIR, ROOT, FairRoot, library paths
\end{lstlisting}

\code{\$VMCWORKDIR} points at the repository root and is used throughout the
macros to locate geometry, parameter, and resource files.

\subsection{Build dependencies}

\begin{itemize}[nosep]
  \item \textbf{ROOT} (with Eve, Geom, RGL, Gui, Physics, Matrix, MathCore, Tree)
  \item \textbf{FairRoot} --- the task framework
  \item \textbf{Eigen3} --- required by the vendored OpenKF (UKF)
  \item \textbf{CATIMA} --- energy-loss model used by the UKF (auto-fetched if absent)
  \item \textbf{GENFIT2} --- required for the genfit fitter (see below)
\end{itemize}

\subsection{Enabling genfit}

genfit support is \emph{conditional} at build time. The top-level
\file{CMakeLists.txt} calls \code{find\_package2(PUBLIC GENFIT2)}, and
\file{AtReconstruction/CMakeLists.txt} only compiles \file{AtGenfit.cxx} and
\file{AtSpacePointMeasurement.cxx} when \code{GENFIT2\_FOUND} is true:

\begin{lstlisting}[style=bash]
if(GENFIT2_FOUND)
  set(SRCS ${SRCS} AtFitter/AtGenfit.cxx AtFitter/AtSpacePointMeasurement.cxx ...)
  set(DEPENDENCIES ${DEPENDENCIES} GENFIT2::genfit2)
endif()
\end{lstlisting}

\begin{warnbox}
To build with genfit you must have GENFIT2 installed and the \code{GENFIT}
environment variable (or CMake variable) pointing at its install prefix --- the
finder expects \code{\$GENFIT/lib/libgenfit2.so} and \code{\$GENFIT/include/}.
If \code{GENFIT2} is not found, the framework still builds, but the genfit fitter
is simply absent and only the UKF is available.
\end{warnbox}

% =====================================================================
\section{Simulation pipeline (16C+p example)}\label{sec:sim}
% =====================================================================

This section reconstructs simulated $^{16}$C$(p,p)$ data. All macros live in
\file{macro/Simulation/ATTPC/16C\_pp/}. Run them from that directory.

\subsection{Stage 1 --- Simulation}

\textbf{Macro:} \file{C16\_pp\_sim.C} \quad
\code{void C16\_pp\_sim(Int\_t nEvents = 100, TString mcEngine = "TGeant4")}

\begin{lstlisting}[style=bash]
root -b -q 'C16_pp_sim(10000)'
\end{lstlisting}

It runs the Geant4 transport for a $^{16}$C beam on a hydrogen target with the
$(p,p)$ reaction generator, in a 28.5\,kG ($2.85$\,T) axial field. Key setup:

\begin{lstlisting}
// 16C beam
AtTPCIonGenerator *ionGen = new AtTPCIonGenerator("Ion", 6,16,0,1, 0,0,2.297, ...);
// (p,p) two-body reaction, CM 20-90 deg
AtTPC2Body *TwoBody = new AtTPC2Body("TwoBody", ..., 40.0, 20.0, 90.0);
// axial field B_z = 28.5 kG
AtConstField *fMagField = new AtConstField();
fMagField->SetField(0., 0., 28.5);
fMagField->SetFieldRegion(-50,50, -50,50, -10,230);
run->SetField(fMagField);
ATTPC->SetGeometryFileName("ATTPC_H300torr.root");
\end{lstlisting}

\textbf{Output:} \file{data/attpcsim.root} (MC truth), \file{data/attpcpar.root}
(parameters), \file{data/geofile\_full.root} (geometry).

\subsection{Stage 2 --- Digitization + PSA + PRA}

\textbf{Macro:} \file{run\_digi\_attpc.C} \quad
\code{void run\_digi\_attpc(int nEvents = 10000, float tCluster = 8.0)}

\begin{lstlisting}[style=bash]
root -b -q 'run_digi_attpc.C(10000, 8.0)'
\end{lstlisting}

\textbf{Input:} \file{data/attpcsim.root}. Task chain (in order):

\begin{lstlisting}
AtClusterizeTask *clusterizer = new AtClusterizeTask();    // MC -> e- clusters
clusterizer->SetPersistence(kFALSE);
fRun->AddTask(clusterizer);

AtPulseTask *pulse = new AtPulseTask(std::make_shared<AtPulse>(mapping)); // -> pads
pulse->SetPersistence(kTRUE);  pulse->SetSaveMCInfo();
fRun->AddTask(pulse);

auto psa = std::make_unique<AtPSAMax>();   psa->SetThreshold(0);  // hits
AtPSAtask *psaTask = new AtPSAtask(std::move(psa));
psaTask->SetPersistence(kTRUE);
fRun->AddTask(psaTask);

AtPRAtask *praTask = new AtPRAtask();      // track candidates (RANSAC/triplet)
praTask->SetTcluster(tCluster);            // time-clustering window (us)
praTask->SetPersistence(kTRUE);
fRun->AddTask(praTask);
\end{lstlisting}

Parameter file \file{\$VMCWORKDIR/parameters/ATTPC.e20009\_sim.par}, pad map
\file{\$VMCWORKDIR/scripts/Lookup20150611.xml}.

\textbf{Output:} \file{data/output\_digi.root} with branches \code{AtEventH}
(hits) and \code{AtPatternEvent} (candidates, ready to fit).

\subsection{Stage 3 --- UKF fit}

\textbf{Macro:} \file{run\_ukf\_only.C} \quad \code{void run\_ukf\_only(int nEvents = 10000)}

\begin{lstlisting}[style=bash]
root -b -q 'run_ukf_only.C(10000)'
\end{lstlisting}

\textbf{Input:} \file{data/output\_digi.root}. The fitter task:

\begin{lstlisting}
double charge = 1.602176634e-19;   // C
double mass   = 938.272;           // MeV/c^2 (proton)

auto eloss = std::make_unique<AtTools::AtELossCATIMA>(3.553e-5);  // g/cm^3, H2 300 torr
eloss->SetProjectile(1, 1, 1);                                    // Z, A, q
eloss->SetMaterial({ std::make_tuple(1,1,1) });                   // hydrogen

auto ukf = std::make_unique<EventFit::AtFitterUKF>(charge, mass, std::move(eloss));
ukf->SetBField(ROOT::Math::XYZVector(0, 0, 2.85));   // +2.85 T : SIMULATION sign
ukf->SetUKFParameters(1e-3, 2.0, 0.0);               // alpha, beta, kappa
ukf->SetMeasurementSigma(2.0);                       // mm
ukf->SetMomentumSigmaFrac(0.5);                      // loose seed
ukf->SetEnableEnergyStraggling(false);
ukf->SetMinClusters(10);
ukf->SetNIterations(1);
ukf->SetZPadPlane(1000.0);                           // digi->lab z conversion

AtFitterTask *fitterTask = new AtFitterTask(std::move(ukf));
fitterTask->SetInputBranch("AtPatternEvent");
fitterTask->SetOutputBranch("AtTrackingEvent");
fitterTask->SetFitMetadataBranch("AtFitMetadata");
fitterTask->SetPersistence(kTRUE);
fRun->AddTask(fitterTask);
\end{lstlisting}

\textbf{Output:} \file{data/output\_ukf\_only.root}, branch \code{AtTrackingEvent}
(fitted $p,\theta,\phi,$ KE) plus \code{AtFitMetadata}.

\begin{notebox}
\textbf{Validated performance} on this channel ($\sim$10k events): efficiency
$\approx 95.6\%$, KE bias $-0.2\%$ with RMS $\approx 5\%$, $\theta$ RMS
$\approx 0.7^\circ$.
\end{notebox}

\subsection{Stage 4 --- Display and analysis}

\begin{lstlisting}[style=bash]
root -l 'run_ukf_display.C(1)'   # interactive: clusters + fit + MC truth (event 1)
root -l 'display_ukf.C(3)'       # WSL-safe static projections (XY/XZ/YZ + 3D)
root -b -q 'run_ukf_perf.C(500)' # timing / performance stats
root -b -q show_kinematics.C     # KE vs theta, resolution plots
\end{lstlisting}

\subsection{One-shot command sequence}

\begin{lstlisting}[style=bash]
cd macro/Simulation/ATTPC/16C_pp/
root -b -q 'C16_pp_sim(10000)'          # 1. simulate
root -b -q 'run_digi_attpc.C(10000,8.0)'# 2. digi + PSA + PRA
root -b -q 'run_ukf_only.C(10000)'      # 3. UKF fit
root -l   'run_ukf_display.C(1)'        # 4. inspect
\end{lstlisting}

\begin{notebox}
There are also combined macros (\file{run\_digi\_ukf.C}, \file{run\_reco\_ukf.C})
that fold digitization and fitting into one process. They are convenient but the
single-process digi+UKF chain has been observed to hang in some configurations;
the staged sequence above is the robust path.
\end{notebox}

% =====================================================================
\section{Experimental pipeline (a1975 example)}\label{sec:exp}
% =====================================================================

This section reconstructs experimental $^{16}$C+$p$ data (experiment a1975) from
raw HDF5. Macros are in \file{macro/Unpack\_HDF5/a1975/UKF/}; run them from that
folder (the ``UKF root''), as data is read by bare name and plots are written to
each subfolder's \file{plots/}. The reusable infrastructure is under
\file{pipeline/}.

\subsection{The staged unpacking macros}

The pipeline is built incrementally so each stage can be validated and cached. All
read raw HDF5 from \file{/mnt/f/a1975/h5/<run>.h5}.

\begin{longtable}{@{}p{4.2cm} p{3.2cm} p{7.2cm}@{}}
\toprule
\textbf{Macro} & \textbf{Output} & \textbf{Tasks (in order)} \\
\midrule
\endhead
\file{unpack\_a1975\_UKF.C}    & \file{<run>\_unpack.root} & Unpack only \\
\file{unpackPSA\_a1975\_UKF.C} & \file{<run>\_psa.root} / \file{\_disp.root} & Unpack + PSA \\
\file{unpackReco\_a1975\_UKF.C}& \file{<run>\_reco.root}   & Unpack + PSA + SC + PRA \\
\file{unpackNFit\_a1975\_UKF.C}& \file{<run>\_ukf\_full.root}& Unpack + PSA + SC + PRA + UKF \\
\bottomrule
\end{longtable}

\code{unpackReco\_a1975\_UKF.C} is the \textbf{standard entry point}: it produces
the \code{AtPatternEvent} that both fitters consume. Its task chain:

\begin{lstlisting}
auto unpacker = std::make_unique<AtHDFUnpacker>(fAtMapPtr);
unpacker->SetInputFileName(inputFile.Data());
unpacker->SetNumberTimestamps(2);
unpacker->SetBaseLineSubtraction(true);
auto unpackTask = new AtUnpackTask(std::move(unpacker));
unpackTask->SetPersistence(persistRaw);
fRun->AddTask(unpackTask);

auto psa = new AtPSAMax();  psa->SetThreshold(20);
AtPSAtask *psaTask = new AtPSAtask(psa);
psaTask->SetPersistence(persistRaw);
psaTask->SetOutputBranch("AtEventH");
fRun->AddTask(psaTask);

auto SCModel = std::make_unique<AtEDistortionModel>();
SCModel->SetCorrectionMaps(zlutFile, radlutFile, tralutFile);   // a1954 LUTs
auto SCTask = new AtSpaceChargeCorrectionTask(std::move(SCModel));
SCTask->SetInputBranch("AtEventH");                              // -> AtEventCorrected
fRun->AddTask(SCTask);

AtPRAtask *praTask = new AtPRAtask();
praTask->SetInputBranch("AtEventCorrected");
praTask->SetOutputBranch("AtPatternEvent");
praTask->SetClusterRadius(15.0);
praTask->SetClusterDistance(7.5);
praTask->SetPersistence(true);
fRun->AddTask(praTask);
\end{lstlisting}

Resource files: geometry \file{ATTPC\_H1bar\_geomanager.root}, map
\file{ANL2023.xml}, parameters \file{ATTPC.a1954.par}, space-charge LUTs in
\file{\$VMCWORKDIR/resources/corrections/a1954/}.

\begin{lstlisting}[style=bash]
# produce the AtPatternEvent once (persistRaw=true keeps raw for displays)
root -b -q 'pipeline/unpackReco_a1975_UKF.C("run_0116", 1000, true)'
\end{lstlisting}

\subsection{Fitter A --- UKF}

\textbf{Macro:} \file{pipeline/fitUKF\_a1975.C}. It reads \file{<run>\_reco.root}
and writes \file{<run>\_ukf<suffix>.root}. Signature (abbreviated):

\begin{lstlisting}
void fitUKF_a1975(TString fileName="run_0116", Long64_t nEvents=-1,
                  TString particles="proton", Int_t bFieldSign=-1,
                  Double_t bFieldMag=2.85, Double_t gasDensity=9.0e-5,
                  TString outSuffix="", TString ioDir="",
                  Double_t measSigma=2.0, Double_t momSigmaFrac=0.3,
                  Int_t nIter=1, Int_t minClusters=10,
                  TString outDir="", Bool_t refTrack=false,
                  Double_t refInflation=4.0);
\end{lstlisting}

Each particle name builds an \code{AtFitterUKF} hypothesis (CATIMA energy loss,
mass and charge from a particle table), and they are combined in an
\code{AtFitterUKFMulti} that keeps the best $\chi^2/\mathrm{ndf}$ per track:

\begin{lstlisting}
auto multi = std::make_unique<EventFit::AtFitterUKFMulti>();
for (each name in particles) {
   auto ukf = MakeHypothesis(name, bFieldSign, bFieldMag, gasDensity,
                             measSigma, momSigmaFrac, nIter, minClusters,
                             refTrack, refInflation);
   multi->AddHypothesis(std::move(ukf), name, /*chargeSign=*/0);
}
AtFitterTask *fitterTask = new AtFitterTask(std::move(multi));
fitterTask->SetInputBranch("AtPatternEvent");
fitterTask->SetOutputBranch("AtTrackingEvent");
fitterTask->SetFitMetadataBranch("AtFitMetadata");
fitterTask->SetPersistence(kTRUE);
\end{lstlisting}

Inside \code{MakeHypothesis} the decisive call for experimental data is the
\textbf{field sign} (see Section~\ref{sec:conventions}):

\begin{lstlisting}
ukf->SetBField(ROOT::Math::XYZVector(0, 0, bFieldSign * bFieldMag)); // -1 for exp
\end{lstlisting}

\begin{lstlisting}[style=bash]
# fast default fit (proton, exp field sign)
root -b -q 'pipeline/fitUKF_a1975.C("run_0116", -1, "proton", -1)'
# resolution-optimized (tighter sigmas), separate output file
root -b -q 'pipeline/fitUKF_a1975.C("run_0116", -1, "proton", -1, 2.85, 9.0e-5, "_opt", "", 0.5, 0.1, 1, 10)'
\end{lstlisting}

\begin{notebox}
\textbf{Resolution tuning.} \code{measSigma=0.5}\,mm with \code{momSigmaFrac=0.1}
gives noticeably better excitation-energy FWHM than the conservative defaults
(\code{2.0}\,mm, \code{0.3}). Use \code{outSuffix} to avoid clobbering files when
scanning parameters.
\end{notebox}

\subsection{Fitter B --- genfit}

\textbf{Macro:} \file{pipeline/fitGenfit\_a1975.C}. Same input
(\file{<run>\_reco.root}), output \file{<run>\_genfit<suffix>.root}. Signature:

\begin{lstlisting}
void fitGenfit_a1975(TString fileName="run_0106", Long64_t nEvents=-1,
                     TString particle="deuteron", Double_t magneticField=2.85,
                     Double_t gasMediumDensity=0.083147, TString ioDir="",
                     TString outSuffix="", TString outDir="",
                     TString elossName="proton_D2_600torr.txt",
                     Int_t minIter=5, Int_t maxIter=20);
\end{lstlisting}

It adds an \code{AtPIDTask} (caches Spyral PID observables) and then the
\code{AtGenfit} fitter wrapped in \code{AtFitterTaskOld}:

\begin{lstlisting}
std::string elossFile = dir + "/resources/energy_loss/" + elossName;
auto fitter = std::make_unique<AtFITTER::AtGenfit>(
    magneticField, 0.00001, 1000.0, elossFile, gasMediumDensity,
    pdg, minIter, maxIter, /*noMatEffects=*/1);
fitter->SetIonName(particle);
fitter->SetMass(mass);              // atomic mass units
fitter->SetAtomicNumber(Z);
fitter->SetNumFitPoints(1.0);
fitter->SetSimulationConvention(0); // experimental data
fitter->EnableMerging(1);
fitter->EnableSingleVertexTrack(1);
fitter->EnableReclustering(1, 15.0, 7.5);

AtFitterTaskOld *fitterTask = new AtFitterTaskOld(std::move(fitter));
fitterTask->SetInputBranch("AtPatternEvent");
fitterTask->SetOutputBranch("AtTrackingEvent");
fitterTask->SetPersistence(kTRUE);
\end{lstlisting}

\begin{lstlisting}[style=bash]
root -b -q 'pipeline/fitGenfit_a1975.C("run_0106", 200, "deuteron")'
\end{lstlisting}

\subsection{Comparing the two fitters}

Because both read the identical \code{AtPatternEvent}, any difference in the
result is purely the fitter:

\begin{lstlisting}[style=bash]
root -b -q 'pipeline/unpackReco_a1975_UKF.C("run_0106", 500, true)'  # once
root -b -q 'pipeline/fitUKF_a1975.C("run_0106", -1, "deuteron", -1, 2.85, 9.0e-5, "", "", 0.5, 0.1, 1, 10)'
root -b -q 'pipeline/fitGenfit_a1975.C("run_0106", -1, "deuteron")'
\end{lstlisting}

\begin{center}
\begin{tabular}{@{}lll@{}}
\toprule
 & \textbf{UKF} (\code{AtFitterUKF}) & \textbf{genfit} (\code{AtGenfit}) \\
\midrule
Algorithm     & Unscented KF + RTS smoother & \code{KalmanFitterRefTrack} \\
Energy loss   & CATIMA (model)              & SRIM table (file) \\
Iterations    & \code{nIter} (1--2)         & \code{minIter}--\code{maxIter} (5--20) \\
Field sign    & \code{bFieldSign} ($-1$ exp)& \code{SetSimulationConvention(0)} \\
Re-clustering & internal / adaptive         & \code{EnableReclustering(\dots)} \\
Task wrapper  & \code{AtFitterTask}         & \code{AtFitterTaskOld} \\
Output        & \code{AtTrackingEvent} + \code{AtFitMetadata} & \code{AtTrackingEvent} \\
\bottomrule
\end{tabular}
\end{center}

\subsection{Event display}

\begin{lstlisting}[style=bash]
# interactive EVE viewer (needs OpenGL/X11; not available under WSL)
root -l 'pipeline/run_eve_UKF.C("run_0116_disp")'
# WSL-safe static PNGs:
root -b -q 'pipeline/event_png_UKF.C("run_0116_disp", -1)'        # busiest event
root -b -q 'pipeline/event_png_batch_UKF.C("run_0116_disp", 0, 24)' # contact sheet
\end{lstlisting}

\begin{warnbox}
Under WSL there is no working OpenGL/Eve --- interactive 3D viewers crash. Use the
static PNG macros (\file{event\_png\_UKF.C}, \file{event\_png\_batch\_UKF.C})
instead. Both require a \file{<run>\_disp.root} produced with \code{persistRaw=true}.
\end{warnbox}

% =====================================================================
\section{Fitter API reference}\label{sec:fitterapi}
% =====================================================================

\subsection{\texttt{AtFitterTask} (the FairRoot wrapper)}

\code{AtFitterTask} owns any \code{EventFit::AtFitter} and drives it per event.

\begin{lstlisting}
AtFitterTask(std::unique_ptr<EventFit::AtFitter> fitter);   // takes ownership

void SetInputBranch(TString);       // default "AtPatternEvent"
void SetOutputBranch(TString);      // default "AtTrackingEvent"
void SetFitMetadataBranch(TString); // enable AtFitMetadata persistence
void SetRawEventBranch(TString);    // optional, fitter access to AtRawEvent
void SetEventBranch(TString);       // optional, fitter access to AtEvent
void SetPersistence(Bool_t = kTRUE);
\end{lstlisting}

It reads a \code{TClonesArray} of \code{AtPatternEvent} and writes
\code{AtTrackingEvent}. (genfit uses the analogous \code{AtFitterTaskOld}.)

\subsection{\texttt{EventFit::AtFitterUKF}}

\textbf{Constructor:}
\begin{lstlisting}
AtFitterUKF(double charge,        // Coulombs (Z * 1.602176634e-19)
            double mass_MeV,      // MeV/c^2
            std::unique_ptr<AtTools::AtELossModel> elossModel);
\end{lstlisting}

\textbf{Most-used setters} (defaults in parentheses):

\begin{longtable}{@{}p{6.4cm} p{8.0cm}@{}}
\toprule
\textbf{Method} & \textbf{Meaning} \\
\midrule
\endhead
\code{SetBField(XYZVector)}            & Magnetic field, Tesla (default $\{0,0,2.85\}$). Sign matters --- see \S\ref{sec:conventions}. \\
\code{SetEField(XYZVector)}            & Electric field, V/m ($\{0,0,0\}$). \\
\code{SetUKFParameters(a,b,k)}         & Unscented transform $\alpha,\beta,\kappa$ ($10^{-3},2,0$). \\
\code{SetMeasurementSigma(mm)}         & Hit position uncertainty (1.0). Lower $\Rightarrow$ better resolution. \\
\code{SetMomentumSigmaFrac(f)}         & Fractional momentum seed uncertainty (0.1). \\
\code{SetMinClusters(n)}               & Skip tracks with fewer clusters (3). \\
\code{SetNIterations(n)}               & Filter passes (1); $>1$ re-seeds with previous result. \\
\code{SetEnableEnergyStraggling(b)}    & Energy straggling in propagator (false). \\
\code{SetMomentumSeed(p\_MeV)}         & Override Brho seed if $>0$. \\
\code{SetZPadPlane(z)}                 & Pad-plane $z$ (mm) for digi$\rightarrow$lab; $z_\mathrm{lab}=z-z_\mathrm{digi}$. \\
\code{SetAdaptiveClustering(b)}        & Re-cluster inside the fit (true); disable for pre-clustered (GNN) input. \\
\code{SetUseArcWalk(b)} / \code{SetArcWalkParams(\dots)} & kNN arc-ordering clustering for centered vertices. \\
\code{SetUseHelixBackExtrap(b)}        & Helix (not linear) back-extrapolation to vertex; needed at high curvature. \\
\code{SetBackExtrapMaxPath(mm)}        & Max back-extrapolation length to vertex (50). \\
\code{SetForceVertexOnBeamAxis(b)}     & Force vertex to $xy=(0,0)$. \\
\code{SetUseRefTrack(b)} / \code{SetUseRefTrackWarmStart(b)} & genfit-style re-linearization around the previous smoothed trajectory (needs \code{nIter}$\ge 2$). \\
\bottomrule
\end{longtable}

\code{AtFitterUKF::Init()} is a no-op (the filter is created lazily on first fit).

\subsection{\texttt{EventFit::AtFitterUKFMulti} (multi-hypothesis)}

Runs several \code{AtFitterUKF} hypotheses on each track and keeps the best
$\chi^2/\mathrm{ndf}$:

\begin{lstlisting}
void AddHypothesis(std::unique_ptr<AtFitterUKF> fitter,
                   std::string name, int chargeSign = 0); // 0 = no PRA-sign filter
void SetChi2Cap(double cap);
std::size_t GetNHypotheses() const;
\end{lstlisting}

\begin{notebox}
Pass \code{chargeSign = 0} for experimental data: the PRA charge sign is
calibrated on the simulation handedness and is unreliable on experiment until
independently validated, so all hypotheses should be evaluated equally.
\end{notebox}

\subsection{\texttt{AtFITTER::AtGenfit}}

\textbf{Constructor:}
\begin{lstlisting}
AtGenfit(Float_t magfield,        // Tesla (stored internally as 10x)
         Float_t minbrho,         // min Brho, Tm
         Float_t maxbrho,         // max Brho, Tm
         std::string eLossFile,   // SRIM-style energy-loss table
         Float_t gasMediumDensity,// g/cm^3
         Int_t pdg   = 2212,      // 2212 p, 1000010020 d, 1000020040 alpha
         Int_t minit = 5,         // min Kalman iterations
         Int_t maxit = 20,        // max Kalman iterations
         Bool_t noMatEffects = kFALSE);
\end{lstlisting}

\textbf{Key setters:} \code{SetIonName}, \code{SetMass} (amu),
\code{SetAtomicNumber}, \code{SetPDGCode}, \code{SetMinIterations} /
\code{SetMaxIterations}, \code{SetMinBrho} / \code{SetMaxBrho},
\code{SetNumFitPoints} (fraction of hits, 0--1), \code{SetSimulationConvention}
(true = simulation $\theta$ convention; false = experiment), \code{EnableMerging},
\code{EnableSingleVertexTrack}, \code{EnableReclustering(on, radius, size)},
\code{SetVerbosityLevel}.

It uses \code{genfit::KalmanFitterRefTrack} (DAF is available but not
instantiated). Measurements come from \code{AtSpacepointMeasurement}, which
converts \code{AtHitCluster} positions mm$\rightarrow$cm and uses a fixed
covariance (diagonal $1,1,1.28$, the larger $z$ term reflecting drift-time
resolution).

\begin{warnbox}
\code{AtGenfit} requires \code{TGeoManager} (geometry) to be loaded before
construction --- it caches material properties for all volumes. The class is
marked \code{[[deprecated]]} but remains the production genfit path.
\end{warnbox}

% =====================================================================
\section{Coordinate conventions and the magnetic-field sign}\label{sec:conventions}
% =====================================================================

\begin{warnbox}
This is the single most common source of unphysical fits. Read it before fitting
\emph{any} experimental data.
\end{warnbox}

Simulation digitization reverses the drift coordinate
($z_\mathrm{digi} = z_\mathrm{pad plane} - z_\mathrm{MC}$, in \code{AtClusterize}),
but experimental data does \emph{not} --- its $z$ comes straight from the drift
time. The two therefore have \textbf{opposite helix handedness}. To make the UKF
helix match the data, flip the sign of $B_z$ for experimental data:

\begin{lstlisting}
// simulation:
ukf->SetBField({0, 0, +2.85});   // bFieldSign = +1
// experiment:
ukf->SetBField({0, 0, -2.85});   // bFieldSign = -1  (the a1975 default)
\end{lstlisting}

Empirically (a1975 run\_0116):

\begin{center}
\begin{tabular}{@{}lccc@{}}
\toprule
$B_z$ & median $\chi^2/\mathrm{ndf}$ & median KE & physical? \\
\midrule
$-2.85$\,T (exp sign) & $\approx 0.04$ & $\approx 7.9$\,MeV & yes \\
$+2.85$\,T (sim sign) & $\approx 3.09$ & $\approx 152$\,MeV & no \\
\bottomrule
\end{tabular}
\end{center}

For genfit, the analogous control is \code{SetSimulationConvention(false)} for
experimental data.

% =====================================================================
\section{Output data model}\label{sec:datamodel}
% =====================================================================

\begin{longtable}{@{}p{3.6cm} p{3.0cm} p{7.4cm}@{}}
\toprule
\textbf{Branch} & \textbf{Produced by} & \textbf{Contents} \\
\midrule
\endhead
\code{AtRawEvent}       & Unpack / Pulse        & Raw pad waveforms (ADC vs time bucket). \\
\code{AtEventH}         & \code{AtPSAtask}      & 3D hits with charge after PSA. \\
\code{AtEventCorrected} & SC correction         & Hits after space-charge distortion correction. \\
\code{AtPatternEvent}   & \code{AtPRAtask}      & Clustered track candidates (fitter input). \\
\code{AtTrackingEvent}  & \code{AtFitterTask}   & \code{AtFittedTrack}: momentum, vertex, kinematics. \\
\code{AtFitMetadata}    & \code{AtFitterTask}   & $\chi^2/\mathrm{ndf}$, iterations, convergence flag, hypothesis. \\
\code{AtPIDEvent}       & \code{AtPIDTask}      & Spyral PID observables (Brho, $\sqrt{dE/dx}$, \dots). \\
\bottomrule
\end{longtable}

% =====================================================================
\section{Quick-reference command cards}\label{sec:quickref}
% =====================================================================

\subsection*{Simulation ($16$C+$p$)}
\begin{lstlisting}[style=bash]
cd macro/Simulation/ATTPC/16C_pp/
root -b -q 'C16_pp_sim(10000)'            # MC truth         -> data/attpcsim.root
root -b -q 'run_digi_attpc.C(10000,8.0)'  # digi+PSA+PRA     -> data/output_digi.root
root -b -q 'run_ukf_only.C(10000)'        # UKF fit          -> data/output_ukf_only.root
root -l   'run_ukf_display.C(1)'          # inspect
\end{lstlisting}

\subsection*{Experiment (a1975), UKF}
\begin{lstlisting}[style=bash]
cd macro/Unpack_HDF5/a1975/UKF/
root -b -q 'pipeline/unpackReco_a1975_UKF.C("run_0116", 1000, true)'  # -> run_0116_reco.root
root -b -q 'pipeline/fitUKF_a1975.C("run_0116", -1, "proton", -1)'    # -> run_0116_ukf.root
root -b -q 'pipeline/event_png_UKF.C("run_0116_disp", -1)'            # static display
\end{lstlisting}

\subsection*{Experiment (a1975), genfit}
\begin{lstlisting}[style=bash]
cd macro/Unpack_HDF5/a1975/UKF/
root -b -q 'pipeline/unpackReco_a1975_UKF.C("run_0106", 500, true)'   # shared input
root -b -q 'pipeline/fitGenfit_a1975.C("run_0106", -1, "deuteron")'   # -> run_0106_genfit.root
\end{lstlisting}

\subsection*{One-shot experiment (no iteration)}
\begin{lstlisting}[style=bash]
root -b -q 'pipeline/unpackNFit_a1975_UKF.C("run_0116", 1000, "proton", -1, 2.85, 9.0e-5)'
\end{lstlisting}

% =====================================================================
\section{Troubleshooting}\label{sec:trouble}
% =====================================================================

\begin{description}[leftmargin=0.5cm,style=nextline]
  \item[Unphysical KE / huge $\chi^2$ on experimental data]
  Almost always the magnetic-field sign. Use \code{bFieldSign = -1} (UKF) or
  \code{SetSimulationConvention(false)} (genfit). The elastic recoil locus should
  sit at a few MeV; see \S\ref{sec:conventions}.
  \item[``Input file not found / run unpackReco first'']
  The fitter macros require \file{<run>\_reco.root}. Run
  \code{unpackReco\_a1975\_UKF.C} once to produce it.
  \item[Event display crashes / hangs (WSL)]
  No OpenGL under WSL. Use \file{event\_png\_UKF.C} /
  \file{event\_png\_batch\_UKF.C} instead of \file{run\_eve\_UKF.C}.
  \item[genfit fitter is missing]
  The build did not find GENFIT2. Set \code{\$GENFIT} and rebuild
  (\S\ref{sec:env}); without it only the UKF is compiled.
  \item[Poor excitation-energy resolution]
  Tighten the UKF measurement model: \code{measSigma}$\approx0.5$\,mm and
  \code{momSigmaFrac}$\approx0.1$ instead of the conservative defaults.
  \item[Fitter settings that helped one channel hurt another]
  UKF/PRA knobs are physics-dependent --- always re-validate per reaction and
  per field setting; do not copy tuning blindly between setups.
\end{description}

\vfill
\begin{center}\small\itshape
Source paths in this manual are relative to the ATTPCROOTv2-OpenKF repository
root. Code excerpts are condensed for readability; consult the macros and the
\file{AtReconstruction/AtFitter/} headers for exact, current signatures.
\end{center}

\end{document}
